% Studies-in-Algebra-and-Group-Theory - Lecture 15: Restriction and Induction of Representations
% Author: S. D. V. Stephens
% Date: 2026-03-10
% Compile with: pdflatex -shell-escape

\documentclass{report}
\input{~/university/preamble.tex}
% preamble-darkmode.tex
% Dark mode extension for lecture notes
% Usage: % preamble-darkmode.tex
% Dark mode extension for lecture notes
% Usage: % preamble-darkmode.tex
% Dark mode extension for lecture notes
% Usage: \input{~/university/preamble-darkmode.tex} AFTER preamble.tex
%
% This provides:
% - Dark background with light text
% - Material Design color palette
% - Ocean blue gradient colors (p1-p9)
% - Enhanced TikZ/PGFPlots for dark backgrounds
% - qq tcolorbox environment
% - tkz-euclide for geometric constructions

% ===========================================
% DARK MODE PACKAGE
% ===========================================
\usepackage{darkmode}
\enabledarkmode

% ===========================================
% ADDITIONAL PACKAGES
% ===========================================
\usepackage{changepage}
\usepackage{ulem}
\usepackage{colortbl}
\usepackage{tkz-euclide}
\usepackage{capt-of}

% ===========================================
% EXTENDED TIKZ LIBRARIES
% ===========================================
\usetikzlibrary{patterns,3d,backgrounds}
\usepgfplotslibrary{groupplots}

% Upgrade pgfplots compatibility
\pgfplotsset{compat=1.18}

% ===========================================
% OCEAN BLUE GRADIENT (p1-p9)
% ===========================================
\definecolor{p1}{HTML}{caf0f8}
\definecolor{p2}{HTML}{ade8f4}
\definecolor{p3}{HTML}{90e0ef}
\definecolor{p4}{HTML}{48cae4}
\definecolor{p5}{HTML}{00b4d8}
\definecolor{p6}{HTML}{0096c7}
\definecolor{p7}{HTML}{0077b6}
\definecolor{p8}{HTML}{023e8a}
\definecolor{p9}{HTML}{03045e}

% Dark background color
\definecolor{pag}{HTML}{293133}

% ===========================================
% MATERIAL DESIGN COLORS
% ===========================================
% Note: myred, myyellow already defined in main preamble
% These are additional/alternate versions
\definecolor{matred}{HTML}{f44336}
\definecolor{matpink}{HTML}{e81e63}
\definecolor{matpurple}{HTML}{9c27b0}
\definecolor{matdeeppurple}{HTML}{673ab7}
\definecolor{matindigo}{HTML}{3f51b5}
\definecolor{matblue}{HTML}{2196f3}
\definecolor{matlightblue}{HTML}{03a9f4}
\definecolor{matcyan}{HTML}{00bcd4}
\definecolor{matteal}{HTML}{009688}
\definecolor{matgreen}{HTML}{4caf50}
\definecolor{matlightgreen}{HTML}{8bc34a}
\definecolor{matlime}{HTML}{cddc39}
\definecolor{matyellow}{HTML}{ffeb3b}
\definecolor{matamber}{HTML}{ffc107}
\definecolor{matorange}{HTML}{ff9800}
\definecolor{matdeeporange}{HTML}{ff5722}
\definecolor{matbrown}{HTML}{795548}
\definecolor{matgray}{HTML}{9e9e9e}
\definecolor{matbluegray}{HTML}{607d8b}

% ===========================================
% COLORLET FOR TIKZ
% ===========================================
\colorlet{xcol}{blue!60!black}

% ===========================================
% DARK MODE TCOLORBOX
% ===========================================
\newtcolorbox{qq}[2][]{%
    boxrule=0.75pt,
    sharp corners,
    colframe=white,
    colback=pag,
    coltext=white,
    #1,
}

% Dark definition box
\newtcolorbox{darkdfn}[2][]{%
    enhanced,
    boxrule=1pt,
    sharp corners,
    colframe=p5,
    colback=pag,
    coltext=white,
    coltitle=white,
    fonttitle=\bfseries,
    title=#2,
    #1,
}

% Dark theorem box
\newtcolorbox{darkthm}[2][]{%
    enhanced,
    boxrule=1pt,
    sharp corners,
    colframe=matpurple,
    colback=pag,
    coltext=white,
    coltitle=white,
    fonttitle=\bfseries,
    title=#2,
    #1,
}

% ===========================================
% TIKZ 3D FIX
% ===========================================
% Small fix for canvas is xy plane at z
% https://tex.stackexchange.com/a/48776/121799
\makeatletter
\tikzoption{canvas is xy plane at z}[]{%
    \def\tikz@plane@origin{\pgfpointxyz{0}{0}{#1}}%
    \def\tikz@plane@x{\pgfpointxyz{1}{0}{#1}}%
    \def\tikz@plane@y{\pgfpointxyz{0}{1}{#1}}%
    \tikz@canvas@is@plane}
\makeatother

% ===========================================
% CONDITIONAL LINE TIKZSET
% ===========================================
\tikzset{
    conditional line/.style 2 args={
        /utils/exec={\pgfmathsetmacro{\xcoordA}{#1}},
        /utils/exec={\pgfmathsetmacro{\xcoordB}{#2}},
        insert path={
            \ifdim\xcoordA pt>\xcoordB pt
            (\xcoordA,0) -- (\xcoordB,0)
            \fi
        },
    },
}

% ===========================================
% PGFPLOTS DARK MODE DEFAULTS
% ===========================================
\pgfplotsset{
    dark axis/.style={
        axis line style={white},
        tick style={white},
        label style={white},
        grid style={pag!70},
        every axis label/.append style={white},
        every tick label/.append style={white},
    },
}

% ===========================================
% TIKZ EXTERNALIZATION (for fast compilation)
% ===========================================
% This creates .md5 cache files for figures
% Requires: pdflatex -shell-escape
%
% To enable, add AFTER loading this file:
%   \enabletikzexternalize
%
\newcommand{\enabletikzexternalize}{%
  \usetikzlibrary{external}%
  \tikzexternalize[prefix=figures/]%
}

% ===========================================
% CONVENIENT SHORTCUTS FOR DARK PLOTS
% ===========================================
\newcommand{\darkplot}[1]{%
    \begin{tikzpicture}
    \begin{axis}[
        dark axis,
        grid=both,
        major grid style={line width=.2pt,draw=pag!70},
        #1
    ]
}


% ===========================================
% QUICK GEOMETRY MACROS (tkz-euclide)
% ===========================================
% Draw a point: \point{name}{x}{y}
\newcommand{\gpoint}[3]{\tkzDefPoint(#2,#3){#1}}

% Draw line through points: \gline{A}{B}
\newcommand{\gline}[2]{\tkzDrawLine[color=p4](#1,#2)}

% Draw circle: \gcircle{center}{radius}
\newcommand{\gcircle}[2]{\tkzDrawCircle[color=p5](#1,#2)}

% Mark angle: \gangle{A}{B}{C}
\newcommand{\gangle}[3]{\tkzMarkAngle[color=p3,size=0.5](#1,#2,#3)}

% ===========================================
% USAGE EXAMPLE
% ===========================================
% In your document:
%
% \begin{tikzpicture}
%   \begin{axis}[dark axis, ...]
%     \addplot[p4, thick] {sin(deg(x))};
%   \end{axis}
% \end{tikzpicture}
%
% Or with qq box:
% \begin{qq}{My Title}
%   Content here in dark box
% \end{qq}
 AFTER preamble.tex
%
% This provides:
% - Dark background with light text
% - Material Design color palette
% - Ocean blue gradient colors (p1-p9)
% - Enhanced TikZ/PGFPlots for dark backgrounds
% - qq tcolorbox environment
% - tkz-euclide for geometric constructions

% ===========================================
% DARK MODE PACKAGE
% ===========================================
\usepackage{darkmode}
\enabledarkmode

% ===========================================
% ADDITIONAL PACKAGES
% ===========================================
\usepackage{changepage}
\usepackage{ulem}
\usepackage{colortbl}
\usepackage{tkz-euclide}
\usepackage{capt-of}

% ===========================================
% EXTENDED TIKZ LIBRARIES
% ===========================================
\usetikzlibrary{patterns,3d,backgrounds}
\usepgfplotslibrary{groupplots}

% Upgrade pgfplots compatibility
\pgfplotsset{compat=1.18}

% ===========================================
% OCEAN BLUE GRADIENT (p1-p9)
% ===========================================
\definecolor{p1}{HTML}{caf0f8}
\definecolor{p2}{HTML}{ade8f4}
\definecolor{p3}{HTML}{90e0ef}
\definecolor{p4}{HTML}{48cae4}
\definecolor{p5}{HTML}{00b4d8}
\definecolor{p6}{HTML}{0096c7}
\definecolor{p7}{HTML}{0077b6}
\definecolor{p8}{HTML}{023e8a}
\definecolor{p9}{HTML}{03045e}

% Dark background color
\definecolor{pag}{HTML}{293133}

% ===========================================
% MATERIAL DESIGN COLORS
% ===========================================
% Note: myred, myyellow already defined in main preamble
% These are additional/alternate versions
\definecolor{matred}{HTML}{f44336}
\definecolor{matpink}{HTML}{e81e63}
\definecolor{matpurple}{HTML}{9c27b0}
\definecolor{matdeeppurple}{HTML}{673ab7}
\definecolor{matindigo}{HTML}{3f51b5}
\definecolor{matblue}{HTML}{2196f3}
\definecolor{matlightblue}{HTML}{03a9f4}
\definecolor{matcyan}{HTML}{00bcd4}
\definecolor{matteal}{HTML}{009688}
\definecolor{matgreen}{HTML}{4caf50}
\definecolor{matlightgreen}{HTML}{8bc34a}
\definecolor{matlime}{HTML}{cddc39}
\definecolor{matyellow}{HTML}{ffeb3b}
\definecolor{matamber}{HTML}{ffc107}
\definecolor{matorange}{HTML}{ff9800}
\definecolor{matdeeporange}{HTML}{ff5722}
\definecolor{matbrown}{HTML}{795548}
\definecolor{matgray}{HTML}{9e9e9e}
\definecolor{matbluegray}{HTML}{607d8b}

% ===========================================
% COLORLET FOR TIKZ
% ===========================================
\colorlet{xcol}{blue!60!black}

% ===========================================
% DARK MODE TCOLORBOX
% ===========================================
\newtcolorbox{qq}[2][]{%
    boxrule=0.75pt,
    sharp corners,
    colframe=white,
    colback=pag,
    coltext=white,
    #1,
}

% Dark definition box
\newtcolorbox{darkdfn}[2][]{%
    enhanced,
    boxrule=1pt,
    sharp corners,
    colframe=p5,
    colback=pag,
    coltext=white,
    coltitle=white,
    fonttitle=\bfseries,
    title=#2,
    #1,
}

% Dark theorem box
\newtcolorbox{darkthm}[2][]{%
    enhanced,
    boxrule=1pt,
    sharp corners,
    colframe=matpurple,
    colback=pag,
    coltext=white,
    coltitle=white,
    fonttitle=\bfseries,
    title=#2,
    #1,
}

% ===========================================
% TIKZ 3D FIX
% ===========================================
% Small fix for canvas is xy plane at z
% https://tex.stackexchange.com/a/48776/121799
\makeatletter
\tikzoption{canvas is xy plane at z}[]{%
    \def\tikz@plane@origin{\pgfpointxyz{0}{0}{#1}}%
    \def\tikz@plane@x{\pgfpointxyz{1}{0}{#1}}%
    \def\tikz@plane@y{\pgfpointxyz{0}{1}{#1}}%
    \tikz@canvas@is@plane}
\makeatother

% ===========================================
% CONDITIONAL LINE TIKZSET
% ===========================================
\tikzset{
    conditional line/.style 2 args={
        /utils/exec={\pgfmathsetmacro{\xcoordA}{#1}},
        /utils/exec={\pgfmathsetmacro{\xcoordB}{#2}},
        insert path={
            \ifdim\xcoordA pt>\xcoordB pt
            (\xcoordA,0) -- (\xcoordB,0)
            \fi
        },
    },
}

% ===========================================
% PGFPLOTS DARK MODE DEFAULTS
% ===========================================
\pgfplotsset{
    dark axis/.style={
        axis line style={white},
        tick style={white},
        label style={white},
        grid style={pag!70},
        every axis label/.append style={white},
        every tick label/.append style={white},
    },
}

% ===========================================
% TIKZ EXTERNALIZATION (for fast compilation)
% ===========================================
% This creates .md5 cache files for figures
% Requires: pdflatex -shell-escape
%
% To enable, add AFTER loading this file:
%   \enabletikzexternalize
%
\newcommand{\enabletikzexternalize}{%
  \usetikzlibrary{external}%
  \tikzexternalize[prefix=figures/]%
}

% ===========================================
% CONVENIENT SHORTCUTS FOR DARK PLOTS
% ===========================================
\newcommand{\darkplot}[1]{%
    \begin{tikzpicture}
    \begin{axis}[
        dark axis,
        grid=both,
        major grid style={line width=.2pt,draw=pag!70},
        #1
    ]
}


% ===========================================
% QUICK GEOMETRY MACROS (tkz-euclide)
% ===========================================
% Draw a point: \point{name}{x}{y}
\newcommand{\gpoint}[3]{\tkzDefPoint(#2,#3){#1}}

% Draw line through points: \gline{A}{B}
\newcommand{\gline}[2]{\tkzDrawLine[color=p4](#1,#2)}

% Draw circle: \gcircle{center}{radius}
\newcommand{\gcircle}[2]{\tkzDrawCircle[color=p5](#1,#2)}

% Mark angle: \gangle{A}{B}{C}
\newcommand{\gangle}[3]{\tkzMarkAngle[color=p3,size=0.5](#1,#2,#3)}

% ===========================================
% USAGE EXAMPLE
% ===========================================
% In your document:
%
% \begin{tikzpicture}
%   \begin{axis}[dark axis, ...]
%     \addplot[p4, thick] {sin(deg(x))};
%   \end{axis}
% \end{tikzpicture}
%
% Or with qq box:
% \begin{qq}{My Title}
%   Content here in dark box
% \end{qq}
 AFTER preamble.tex
%
% This provides:
% - Dark background with light text
% - Material Design color palette
% - Ocean blue gradient colors (p1-p9)
% - Enhanced TikZ/PGFPlots for dark backgrounds
% - qq tcolorbox environment
% - tkz-euclide for geometric constructions

% ===========================================
% DARK MODE PACKAGE
% ===========================================
\usepackage{darkmode}
\enabledarkmode

% ===========================================
% ADDITIONAL PACKAGES
% ===========================================
\usepackage{changepage}
\usepackage{ulem}
\usepackage{colortbl}
\usepackage{tkz-euclide}
\usepackage{capt-of}

% ===========================================
% EXTENDED TIKZ LIBRARIES
% ===========================================
\usetikzlibrary{patterns,3d,backgrounds}
\usepgfplotslibrary{groupplots}

% Upgrade pgfplots compatibility
\pgfplotsset{compat=1.18}

% ===========================================
% OCEAN BLUE GRADIENT (p1-p9)
% ===========================================
\definecolor{p1}{HTML}{caf0f8}
\definecolor{p2}{HTML}{ade8f4}
\definecolor{p3}{HTML}{90e0ef}
\definecolor{p4}{HTML}{48cae4}
\definecolor{p5}{HTML}{00b4d8}
\definecolor{p6}{HTML}{0096c7}
\definecolor{p7}{HTML}{0077b6}
\definecolor{p8}{HTML}{023e8a}
\definecolor{p9}{HTML}{03045e}

% Dark background color
\definecolor{pag}{HTML}{293133}

% ===========================================
% MATERIAL DESIGN COLORS
% ===========================================
% Note: myred, myyellow already defined in main preamble
% These are additional/alternate versions
\definecolor{matred}{HTML}{f44336}
\definecolor{matpink}{HTML}{e81e63}
\definecolor{matpurple}{HTML}{9c27b0}
\definecolor{matdeeppurple}{HTML}{673ab7}
\definecolor{matindigo}{HTML}{3f51b5}
\definecolor{matblue}{HTML}{2196f3}
\definecolor{matlightblue}{HTML}{03a9f4}
\definecolor{matcyan}{HTML}{00bcd4}
\definecolor{matteal}{HTML}{009688}
\definecolor{matgreen}{HTML}{4caf50}
\definecolor{matlightgreen}{HTML}{8bc34a}
\definecolor{matlime}{HTML}{cddc39}
\definecolor{matyellow}{HTML}{ffeb3b}
\definecolor{matamber}{HTML}{ffc107}
\definecolor{matorange}{HTML}{ff9800}
\definecolor{matdeeporange}{HTML}{ff5722}
\definecolor{matbrown}{HTML}{795548}
\definecolor{matgray}{HTML}{9e9e9e}
\definecolor{matbluegray}{HTML}{607d8b}

% ===========================================
% COLORLET FOR TIKZ
% ===========================================
\colorlet{xcol}{blue!60!black}

% ===========================================
% DARK MODE TCOLORBOX
% ===========================================
\newtcolorbox{qq}[2][]{%
    boxrule=0.75pt,
    sharp corners,
    colframe=white,
    colback=pag,
    coltext=white,
    #1,
}

% Dark definition box
\newtcolorbox{darkdfn}[2][]{%
    enhanced,
    boxrule=1pt,
    sharp corners,
    colframe=p5,
    colback=pag,
    coltext=white,
    coltitle=white,
    fonttitle=\bfseries,
    title=#2,
    #1,
}

% Dark theorem box
\newtcolorbox{darkthm}[2][]{%
    enhanced,
    boxrule=1pt,
    sharp corners,
    colframe=matpurple,
    colback=pag,
    coltext=white,
    coltitle=white,
    fonttitle=\bfseries,
    title=#2,
    #1,
}

% ===========================================
% TIKZ 3D FIX
% ===========================================
% Small fix for canvas is xy plane at z
% https://tex.stackexchange.com/a/48776/121799
\makeatletter
\tikzoption{canvas is xy plane at z}[]{%
    \def\tikz@plane@origin{\pgfpointxyz{0}{0}{#1}}%
    \def\tikz@plane@x{\pgfpointxyz{1}{0}{#1}}%
    \def\tikz@plane@y{\pgfpointxyz{0}{1}{#1}}%
    \tikz@canvas@is@plane}
\makeatother

% ===========================================
% CONDITIONAL LINE TIKZSET
% ===========================================
\tikzset{
    conditional line/.style 2 args={
        /utils/exec={\pgfmathsetmacro{\xcoordA}{#1}},
        /utils/exec={\pgfmathsetmacro{\xcoordB}{#2}},
        insert path={
            \ifdim\xcoordA pt>\xcoordB pt
            (\xcoordA,0) -- (\xcoordB,0)
            \fi
        },
    },
}

% ===========================================
% PGFPLOTS DARK MODE DEFAULTS
% ===========================================
\pgfplotsset{
    dark axis/.style={
        axis line style={white},
        tick style={white},
        label style={white},
        grid style={pag!70},
        every axis label/.append style={white},
        every tick label/.append style={white},
    },
}

% ===========================================
% TIKZ EXTERNALIZATION (for fast compilation)
% ===========================================
% This creates .md5 cache files for figures
% Requires: pdflatex -shell-escape
%
% To enable, add AFTER loading this file:
%   \enabletikzexternalize
%
\newcommand{\enabletikzexternalize}{%
  \usetikzlibrary{external}%
  \tikzexternalize[prefix=figures/]%
}

% ===========================================
% CONVENIENT SHORTCUTS FOR DARK PLOTS
% ===========================================
\newcommand{\darkplot}[1]{%
    \begin{tikzpicture}
    \begin{axis}[
        dark axis,
        grid=both,
        major grid style={line width=.2pt,draw=pag!70},
        #1
    ]
}


% ===========================================
% QUICK GEOMETRY MACROS (tkz-euclide)
% ===========================================
% Draw a point: \point{name}{x}{y}
\newcommand{\gpoint}[3]{\tkzDefPoint(#2,#3){#1}}

% Draw line through points: \gline{A}{B}
\newcommand{\gline}[2]{\tkzDrawLine[color=p4](#1,#2)}

% Draw circle: \gcircle{center}{radius}
\newcommand{\gcircle}[2]{\tkzDrawCircle[color=p5](#1,#2)}

% Mark angle: \gangle{A}{B}{C}
\newcommand{\gangle}[3]{\tkzMarkAngle[color=p3,size=0.5](#1,#2,#3)}

% ===========================================
% USAGE EXAMPLE
% ===========================================
% In your document:
%
% \begin{tikzpicture}
%   \begin{axis}[dark axis, ...]
%     \addplot[p4, thick] {sin(deg(x))};
%   \end{axis}
% \end{tikzpicture}
%
% Or with qq box:
% \begin{qq}{My Title}
%   Content here in dark box
% \end{qq}


\course{Studies-in-Algebra-and-Group-Theory}

\begin{document}

\lecture{15}{Restriction and Induction of Representations}

Throughout, \(G\) is a finite group and \(H \leq G\) a subgroup.

\section{Restriction}

\dfn{Restriction}{Given a representation \(V\) of \(G\), the restriction of \(V\) to \(H\), denoted \(\operatorname{Res}^G_H V\), is \(V\) viewed as a representation of \(H\) (forget the action of elements outside \(H\)).}

This defines a functor \(\operatorname{Res}^G_H : \operatorname{Rep}(G) \to \operatorname{Rep}(H)\), where the categories have objects = finite-dimensional complex representations and morphisms = \(G\)-equivariant (resp.\ \(H\)-equivariant) linear maps.

An irreducible representation of \(G\) need not remain irreducible upon restriction to \(H\). The problem we address next is: can we go in the opposite direction?


\section{Induced Representations}

Suppose \(V\) is a representation of \(G\) and \(W \subset V\) is a subspace that is invariant under \(H\) (so \(W\) is a subrepresentation of \(\operatorname{Res}^G_H V\)), but not necessarily under all of \(G\). For \(g \in G\), the subspace \(gW \subset V\) depends only on the coset \(gH\), and \(gW\) is a representation of the conjugate subgroup \(gHg^{-1}\) via \(H \xrightarrow{\rho} \GL(W)\), \(gHg^{-1} \xrightarrow{\sim} \GL(gW)\) (conjugation by \(g\)).

The simplest scenario is when
\[V = \bigoplus_{\sigma \in G/H} \sigma W\]
as a direct sum of vector spaces, where \(\sigma\) ranges over a set of coset representatives. In this case, \(\dim V = [G:H] \cdot \dim W\), and the entire representation \(V\) is completely determined by the \(H\)-representation \(W\).

\dfn{Induced representation}{A representation \(V\) of \(G\), together with a subspace \(W \subset V\) invariant under \(H\) (i.e.\ \(W\) is a subrepresentation of \(\operatorname{Res}^G_H V\)), is said to be induced by \(W \in \operatorname{Rep}(H)\) if
\[V = \bigoplus_{\sigma \in G/H} \sigma W.\]
We write \(V = \operatorname{Ind}^G_H W\). Concretely: fix coset representatives \(\sigma_1, \ldots, \sigma_k\) of \(G/H\) (where \(k = [G:H]\)). Then every \(v \in V\) can be written uniquely as \(v = \sigma_1 w_1 + \cdots + \sigma_k w_k\) with \(w_i \in W\), and \(G\) acts by permuting the summands: if \(g\sigma_i = \sigma_j h\) for some \(h \in H\), then \(g(\sigma_i w) = \sigma_j(hw)\).
}

\thm{Existence and uniqueness of induced representations}{Given a representation \(W\) of \(H\), the induced representation \(V = \operatorname{Ind}^G_H W\) exists and is unique up to isomorphism of \(G\)-representations.}

\pf{Proof}{\emph{Uniqueness.} Suppose \(V\) is a representation of \(G\) and \(W \subset V\) is \(H\)-invariant with \(V = \bigoplus_{i=1}^k \sigma_i W\). For any \(g \in G\), there is a unique \(j\) such that \(g\sigma_i \in \sigma_j H\), i.e.\ \(h = \sigma_j^{-1} g \sigma_i \in H\), and necessarily \(g(\sigma_i w) = \sigma_j(hw)\). This determines the \(G\)-action on \(V\) uniquely from the \(H\)-action on \(W\).

\emph{Existence.} Build \(V = \bigoplus_{i=1}^k \sigma_i W\) where the \(\sigma_i\) are formal symbols (i.e.\ \(V\) consists of \(k\) copies of \(W\) as a vector space). Define \(g \in G\) to act by: given \(g\sigma_i = \sigma_j h\), set \(g(\sigma_i w) = \sigma_j(hw)\). This is well-defined (independent of the choice of coset representatives) and defines a group homomorphism \(G \to \GL(V)\).
}

\subsection{Examples}

\ex{Permutation representation from induction}{The permutation representation of \(G\) on the coset space \(G/H\) equals \(\operatorname{Ind}^G_H U\), where \(U\) is the trivial representation of \(H\). Indeed, the permutation representation has basis \(\{e_{\sigma H}\}_{\sigma \in G/H}\), the element \(e_H\) spans a 1-dimensional subspace fixed by \(H\), and \(gW = \operatorname{span}(e_{gH})\), so \(V = \bigoplus_{gH \in G/H} gW\).
}

\ex{Regular representation from induction}{The regular representation of \(G\) is \(\operatorname{Ind}^G_H R_H\), where \(R_H\) is the regular representation of \(H\). Here \(W = \operatorname{span}\{e_h : h \in H\} \subset V = \operatorname{span}\{e_g : g \in G\}\).
}

\subsection{Properties of Induction}

\mprop{Basic properties}{Let \(H \leq G\), and let \(W, W'\) be representations of \(H\) and \(U\) a representation of \(G\).
\begin{enumerate}[(i)]
    \item \(\operatorname{Ind}^G_H(W \oplus W') \cong \operatorname{Ind}^G_H(W) \oplus \operatorname{Ind}^G_H(W')\).
    \item \(\operatorname{Ind}^G_H(\operatorname{Res}^G_H(U) \otimes W) \cong U \otimes \operatorname{Ind}^G_H(W)\).
    \item In particular, \(\operatorname{Ind}^G_H(\operatorname{Res}^G_H(U)) \cong U \otimes \operatorname{Ind}^G_H(U_H)\), where \(U_H\) is the trivial representation of \(H\). The right side is \(U\) tensored with the permutation representation on \(G/H\).
\end{enumerate}
}

\pf{Proof of (ii)}{Write \(\operatorname{Ind}^G_H(W) = \bigoplus_{\sigma \in G/H} \sigma W\). Then
\[U \otimes \operatorname{Ind}^G_H(W) = \bigoplus_{\sigma \in G/H} (U \otimes \sigma W) = \bigoplus_{\sigma \in G/H} \sigma(U \otimes W),\]
where \(U \otimes W \subset U \otimes \operatorname{Ind}^G_H(W)\) is invariant under \(H\), and as an \(H\)-representation it is \(\operatorname{Res}^G_H(U) \otimes W\).
}

\nt{In general, \(\operatorname{Ind}(W \otimes W') \not\cong \operatorname{Ind}(W) \otimes \operatorname{Ind}(W')\).}

\subsection{Character of an Induced Representation}

\mprop{Character formula for induction}{Let \(W\) be a representation of \(H \leq G\). Then
\[\chi_{\operatorname{Ind}^G_H W}(g) = \frac{1}{|H|} \sum_{\substack{s \in G \\ s^{-1}gs \in H}} \chi_W(s^{-1}gs).\]
}

\pf{Proof}{Choose coset representatives \(\sigma_1, \ldots, \sigma_k\) of \(G/H\). The element \(g \in G\) maps \(\sigma_i W\) to \(\sigma_j W\) where \(j\) is determined by \(g\sigma_i \in \sigma_j H\). If \(i \neq j\), this contributes nothing to the trace. If \(i = j\), then \(h = \sigma_i^{-1} g \sigma_i \in H\) and \(g\) acts on \(\sigma_i W\) as \(\sigma_i h \sigma_i^{-1}\) composed with the identification \(\sigma_i W \cong W\), so the trace on \(\sigma_i W\) is \(\Tr(h|_W) = \chi_W(\sigma_i^{-1} g \sigma_i)\). Summing over all \(\sigma_i\):
\[\chi_{\operatorname{Ind} W}(g) = \sum_{\substack{i \\ \sigma_i^{-1} g \sigma_i \in H}} \chi_W(\sigma_i^{-1} g \sigma_i).\]
Since each coset \(\sigma_i H\) has \(|H|\) elements, and the condition \(\sigma_i^{-1} g \sigma_i \in H\) is the same for any representative of the coset, we can rewrite this as a sum over all \(s \in G\) with the normalization factor \(1/|H|\).
}


\section{Frobenius Reciprocity}

\thm{Frobenius reciprocity}{Let \(U\) be a representation of \(G\) and \(W\) a representation of \(H \leq G\). Then every \(H\)-equivariant map \(\varphi : W \to \operatorname{Res}^G_H(U)\) extends uniquely to a \(G\)-equivariant map \(\tilde{\varphi} : \operatorname{Ind}^G_H(W) \to U\). That is,
\[\Hom_H(W, \operatorname{Res}^G_H U) \cong \Hom_G(\operatorname{Ind}^G_H W, U).\]
}

\pf{Proof}{Choose coset representatives \(\sigma_1, \ldots, \sigma_k \in G\) and write \(V = \operatorname{Ind}^G_H W = \bigoplus_{i=1}^k \sigma_i W\).

\emph{Uniqueness.} Given \(\varphi : W \to \operatorname{Res}(U)\), if \(\tilde{\varphi} : V \to U\) is \(G\)-equivariant with \(\tilde{\varphi}|_W = \varphi\), then for any \(\sigma_i w \in \sigma_i W\):
\[\tilde{\varphi}(\sigma_i w) = \sigma_i \cdot \tilde{\varphi}(w) = \sigma_i \cdot \varphi(w),\]
where the last \(\sigma_i\) acts via the \(G\)-action on \(U\). This determines \(\tilde{\varphi}\) on each summand, hence on all of \(V\).

\emph{Existence and \(G\)-equivariance.} Define \(\tilde{\varphi}(\sigma_i w) = \sigma_i \cdot \varphi(w)\). To verify \(G\)-equivariance: let \(g \in G\) act on \(V\) by sending \(\sigma_i W\) to \(\sigma_j W\) where \(g\sigma_i = \sigma_j h\) for \(h \in H\). For \(\sigma_i w \in \sigma_i W\):
\[\tilde{\varphi}(g(\sigma_i w)) = \tilde{\varphi}(\sigma_j(hw)) = \sigma_j \cdot \varphi(hw) = \sigma_j \cdot h \cdot \varphi(w) = g \cdot \sigma_i \cdot \varphi(w) = g \cdot \tilde{\varphi}(\sigma_i w),\]
using \(H\)-equivariance of \(\varphi\) in the third equality and \(\sigma_j h = g\sigma_i\) in the fourth. So \(\tilde{\varphi}\) is \(G\)-equivariant.

Conversely, given \(\tilde{\varphi} \in \Hom_G(V, U)\), its restriction to \(W \subset V\) is \(H\)-equivariant. These two maps are inverse to each other.

Comparing dimensions of the Hom spaces gives:
}

\cor{Frobenius reciprocity (character form)}{For any representation \(U\) of \(G\) and \(W\) of \(H\),
\[\langle \chi_{\operatorname{Ind}^G_H W}, \chi_U \rangle_G = \langle \chi_W, \chi_{\operatorname{Res}^G_H U} \rangle_H.\]
}

\nt{Thus, if \(U\) is an irreducible representation of \(G\) and \(W\) an irreducible representation of \(H\), the number of times \(W\) appears in \(\operatorname{Res}^G_H(U)\) equals the number of times \(U\) appears in \(\operatorname{Ind}^G_H(W)\).}


\section{Example: \(S_4 \supset S_3\)}

\ex{Restriction and induction between \(S_4\) and \(S_3\)}{Let \(G = S_4\) with irreducibles \(U_4, U_4', V_4, V_4', W\) (dimensions 1, 1, 3, 3, 2), and \(H = S_3\) with irreducibles \(U_3, U_3', V_3\) (dimensions 1, 1, 2). We compute restrictions using character tables (restricting the character of a representation of \(S_4\) to elements of \(S_3 \hookrightarrow S_4\)):
\begin{itemize}
    \item \(\operatorname{Res}(U_4) = U_3\) and \(\operatorname{Res}(U_4') = U_3'\) (trivial and sign restrict to trivial and sign).
    \item \(\operatorname{Res}(V_4) = V_3 \oplus U_3\). Indeed, the permutation representation \(\CC^4\) restricts as \(\CC^3 \oplus \CC\) (the 4th coordinate is fixed by \(S_3\)), so \(\operatorname{Res}(V_4 \oplus U_4) = V_3 \oplus U_3 \oplus U_3\).
    \item \(\operatorname{Res}(V_4') = V_3 \oplus U_3'\) (since \(V_4' = V_4 \otimes U_4'\) and \(V_3 \otimes U_3' \cong V_3\)).
    \item \(\operatorname{Res}(W) = V_3\). The representation \(W\) factors through \(S_4/\{e, (12)(34), (13)(24), (14)(23)\} \cong S_3\), and \(W\) as a representation of this quotient is the standard representation \(V_3\).
\end{itemize}
By Frobenius reciprocity, we can read off the induced representations:
\begin{itemize}
    \item \(\operatorname{Ind}^{S_4}_{S_3}(V_3)\) contains \(V_4, V_4', W\) each with multiplicity 1, giving \(\operatorname{Ind}(V_3) = V_4 \oplus V_4' \oplus W\) (dim \(= 3 + 3 + 2 = 8 = [S_4:S_3] \cdot 2\)).
    \item \(\operatorname{Ind}^{S_4}_{S_3}(U_3) = U_4 \oplus V_4\) (dim \(= 1 + 3 = 4 = [S_4:S_3] \cdot 1\)).
    \item \(\operatorname{Ind}^{S_4}_{S_3}(U_3') = U_4' \oplus V_4'\) (dim \(= 1 + 3 = 4\)).
\end{itemize}
}



\mer{Exercises}{
\begin{enumerate}[(1)]
    \item Let \(H = \{e\} \leq G\). Show that \(\operatorname{Ind}^G_{\{e\}} W \cong W \otimes R\) where \(R\) is the regular representation of \(G\) and \(W\) is viewed as a vector space with trivial \(G\)-action on the second factor.

    \item Let \(H = A_n \leq S_n\) for \(n \geq 2\). Show that \(\operatorname{Ind}^{S_n}_{A_n}(U)\), where \(U\) is the trivial representation of \(A_n\), is isomorphic to \(U_{S_n} \oplus U'_{S_n}\) (the trivial plus the sign representation of \(S_n\)).

    \item Using Frobenius reciprocity and the character table of \(S_5\), compute \(\operatorname{Ind}^{S_5}_{S_4}(V_4)\) (where \(V_4\) is the standard representation of \(S_4\), embedded as the subgroup fixing 5).

    \item Prove that the character formula for induction can alternatively be written as \(\chi_{\operatorname{Ind} W}(g) = \sum_{i : \sigma_i^{-1} g \sigma_i \in H} \chi_W(\sigma_i^{-1} g \sigma_i)\) where \(\sigma_1, \ldots, \sigma_k\) are coset representatives.

    \item (Transitivity of induction) If \(K \leq H \leq G\), show \(\operatorname{Ind}^G_H \operatorname{Ind}^H_K W \cong \operatorname{Ind}^G_K W\).
\end{enumerate}
}

\begin{solution}
\textbf{(1)} \(\operatorname{Ind}^G_{\{e\}} W = \bigoplus_{g \in G} gW\), a direct sum of \(|G|\) copies of \(W\), with \(G\) permuting the summands. This is exactly \(W \otimes R\) where the \(G\)-action is on the \(R\) factor.

\textbf{(2)} \(\operatorname{Ind}^{S_n}_{A_n}(U) = \bigoplus_{\sigma \in S_n/A_n} \sigma U\). Since \([S_n:A_n] = 2\), this is 2-dimensional with basis \(e_1, e_2\) indexed by the two cosets. Even permutations fix both; odd permutations swap them. The invariant subspace \(\operatorname{span}(e_1 + e_2) \cong U_{S_n}\), and the anti-invariant \(\operatorname{span}(e_1 - e_2) \cong U'_{S_n}\).

\textbf{(3)} First compute \(\operatorname{Res}^{S_5}_{S_4}\) of each irreducible of \(S_5\) and find the multiplicity of \(V_4\). From Frobenius reciprocity, the multiplicity of an \(S_5\)-irreducible \(X\) in \(\operatorname{Ind}(V_4)\) equals the multiplicity of \(V_4\) in \(\operatorname{Res}(X)\). Computing (by restricting characters to \(S_4\)): \(\operatorname{Ind}^{S_5}_{S_4}(V_4) = V \oplus V' \oplus W \oplus W' \oplus \bigwedge^2 V\) (\(\dim = 4+4+5+5+6 = 24 ... \) but \([S_5:S_4] \cdot 3 = 15\)). Let me recompute: \([S_5:S_4] = 5\) and \(\dim V_4 = 3\), so \(\dim \operatorname{Ind}(V_4) = 15\). The correct decomposition requires careful character computation.

\textbf{(4)} The sum \(\sum_{i : \sigma_i^{-1}g\sigma_i \in H} \chi_W(\sigma_i^{-1}g\sigma_i)\) counts each valid conjugate exactly once per coset. In the formula \(\frac{1}{|H|}\sum_{s \in G, s^{-1}gs \in H}\), replacing \(s\) by \(\sigma_i h\) for \(h \in H\): the condition \(s^{-1}gs \in H\) becomes \(h^{-1}(\sigma_i^{-1}g\sigma_i)h \in H\), which holds iff \(\sigma_i^{-1}g\sigma_i \in H\). The \(|H|\) values of \(h\) give \(\chi_W(h^{-1}(\sigma_i^{-1}g\sigma_i)h) = \chi_W(\sigma_i^{-1}g\sigma_i)\) (class function), contributing \(|H|\) times the same value. Dividing by \(|H|\) recovers the coset-representative sum.

\textbf{(5)} Write coset representatives for \(G/K\) by combining those for \(G/H\) and \(H/K\): if \(\{\sigma_i\}\) represents \(G/H\) and \(\{\tau_j\}\) represents \(H/K\), then \(\{\sigma_i \tau_j\}\) represents \(G/K\). The direct sum decompositions are compatible: \(\operatorname{Ind}^G_K W = \bigoplus_{i,j} \sigma_i \tau_j W = \bigoplus_i \sigma_i (\bigoplus_j \tau_j W) = \bigoplus_i \sigma_i (\operatorname{Ind}^H_K W) = \operatorname{Ind}^G_H(\operatorname{Ind}^H_K W)\).
\end{solution}

\end{document}
